\scanpage{36}
aux homomorphismes
\diagram{\begin{tikzcd}
T \arrow[r] & A^{\unsure{\natural|\sharp}} \arrow[r] & A
\end{tikzcd}},
\margin{$\unsure{=}$}
nous montre alors aussitôt que $\uline{F}(A^\natural)$ et $\uline{F}(T)$ s'identifient respectivement à
\[
(13.4.10)\qquad\uline{F}(A^\natural)\simeq\uline{F}^{f}\overset{\mathrm{dfn}}=\uline{F}\cap\mathcal E^{f},\quad\uline{F}(T)=\uline{F}^{t}=\uline{F}\cap\mathcal E^{t}=0,
\]
le $=0$ provenant du fait immédiat que pour un tore $T$, on a $\uline{F}(T)=0$,
la suite exacte du type (13.4.3) devenant ici
\diagram{\begin{tikzcd}
0 \arrow[r] & 0 \arrow[r] & \mathcal E(T) \arrow[r] & \uline{\operatorname{Lie}}'(T) \arrow[r] & 0
\end{tikzcd}}.
On voit donc en même temps que l'homomorphisme composé
\diagram{\begin{tikzcd}
\mathcal E^{t} \arrow[r] & \mathcal E \arrow[r] & \uline{E}
\end{tikzcd}}
\margin{$\unsure{-}$}
induit un isomorphisme

\margin{495}
\figure{\begin{tikzpicture}\draw (0,0)--(.36,-.03);\end{tikzpicture}}
\margin{$\unsure{0=}$}
\fbox{495}
\mspage{35}
